\section{Serangan \& Insiden Siber pada Sistem Kesehatan}

\subsection{Definisi Serangan Siber}
Serangan siber (\textit{cyber attack}) dapat didefinisikan sebagai setiap upaya yang disengaja untuk mengeksploitasi kerentanan pada sistem komputer, jaringan, atau infrastruktur digital. Tujuan dari serangan ini beragam, mulai dari mencuri, mengubah, menonaktifkan, atau menghancurkan data. Menurut \textit{National Institute of Standards and Technology} (NIST), serangan siber merupakan tindakan yang dilakukan untuk merusak integritas, kerahasiaan, atau ketersediaan informasi \cite{13}. Serangan ini dapat dilancarkan oleh berbagai aktor, dari peretas individu hingga organisasi yang disponsori negara, dengan menggunakan berbagai teknik dan alat.

\subsection{Definisi Insiden Siber}
Insiden siber (\textit{cyber incident}) adalah peristiwa yang terjadi sebagai akibat dari serangan siber yang berhasil, yang berpotensi atau secara aktual membahayakan sistem informasi. ISO/IEC 27035 mendefinisikan insiden keamanan informasi sebagai satu atau serangkaian kejadian keamanan informasi yang tidak diinginkan atau tidak terduga yang memiliki probabilitas signifikan untuk mengkompromikan operasi bisnis dan mengancam keamanan informasi \cite{18}. Sebuah insiden menandakan bahwa pertahanan keamanan telah ditembus dan sistem mungkin telah dikompromikan. Oleh karena itu, respons yang cepat dan efektif sangat penting untuk meminimalkan dampak.

\subsection{Hubungan antara Serangan dan Insiden Siber}
Serangan dan insiden siber memiliki hubungan sebab-akibat dengan serangan adalah 'sebab', sedangkan insiden adalah 'akibat'. Sebuah serangan yang berhasil menembus sistem pertahanan akan memicu terjadinya sebuah insiden. Sebagai contoh, dalam konteks sistem RME, sebuah serangan \textit{phishing} (sebab) yang berhasil mengelabui seorang staf untuk mengklik tautan berbahaya dapat menyebabkan insiden berupa instalasi \textit{malware} atau akses tidak sah ke data pasien (akibat).